\problemname{Laviner}
\noindent
Alexander och Joshua är på en skidresa tillsammans, och Alexander planerar att göra ``lite'' bus idag. 
Tillsammans har de åkt till ett stort skidsystem med $N$ skidbyar som är sammankopplat med $N-1$ enkelriktade skidbackar. 
Skidby $1$ befinner sig t upp i skidsystemet, och det är möjligt att nå alla andra skidbyar från skidby $1$ genom att endast åka ner för skidbackar.
Alexander får åka upp och ställa sig vid en valfri skidby och skrika. Det skapas då en lavin som utgår från skidbyn och fortsätter längs med alla riktade skidbackar. 

Men Joshua vill förhindra Alexanders bus, och kan placera ut $K$ murar på $K$ olika skidbyar. 
Muren skyddar den specifika skidbyn från att bli påverkad av lavinen och skriket, och skyddar på sätt alla andra skidbyar under den skidbyn. 
Joshua är en skarpsinnig grabb, och kommer att placera murarna på ett sätt som minimerar antal antalet påverkade skidbyar oavsett vilken skidby som Alexbörjar skrika ifrån.
Notera att Alexander kan inte heller starta en lavin och skrika från en skidby där Joshua har murat in.

Vad är det maximala antalet skidbyar som Alexander kan drabba trots att Joshua försöker stoppa honom? 

\section*{Indata}
Den första raden i indatan innehåller två heltal, $N,K$ ($1 \le K < N \le 10^5$), antalet skidbyar i skidsystemet, samt antalet murar Joshua kan sätta ut.

Därefter följer en rad med $N-1$ heltal $a_i$, där för varje $i$ ($1 \leq i \leq N-1$) representerar $a_i$ 
att det finns en enkelriktad skidbacke där man kan åka från skidby $a_i$ ($1 \le a_i \le i$) till skidby $i+1$. 


\section*{Utdata}
Skriv ut ett heltal, maximala antalet skidbyar som Alexander drabba genom att skrika från en nod en gång, efter att Joshua har placerat ut sina $K$ murar.

\section*{Poängsättning}
Din lösning kommer att testas på en mängd testfallsgrupper.
För att få poäng för en grupp så måste du klara alla testfall i gruppen.

\noindent
\begin{tabular}{| l | l | p{12cm} |}
  \hline
  \textbf{Grupp} & \textbf{Poäng} & \textbf{Gränser} \\ \hline
  $1$    & $10$       & $a_i = i$ för alla $(1 \le i \le N-1)$. \\ \hline
  $2$    & $10$       & $K = 1, N \le 1000$ \\ \hline
  $3$    & $10$       & $N \le 20$ \\ \hline
  $4$    & $15$       & $K = 1$ \\ \hline
  $5$    & $25$       & $N \leq 1000$ \\ \hline
  $6$    & $30$       & Inga ytterligare begränsningar. \\ \hline
\end{tabular}

%\section*{Exempelfall}

